\documentclass[12pt]{article}
\usepackage[utf8]{inputenc}
\usepackage{amsmath, amssymb, amsthm}
\usepackage{geometry}

% Set margins to match a standard homework document
\geometry{a4paper, margin=1in}

\title{Workshop 7}
\parskip=1ex
\author{Math 292 Spring 2026 \\ Rutgers University}
\date{}

\begin{document}

\maketitle

\noindent{}\textbf{Exercise 1:}
Let $[a,b]$ be a bounded interval.
Suppose $p(x)$ is strictly positive and continuously differentiable on $[a,b]$.
Consider
\begin{align*}
    \mathcal{L}u(x) = (p(x)u'(x))'
\end{align*}
and the eigenfunctions $\mathcal{L}u(x) = \lambda u(x)$ with Neumann boundary conditions
\begin{align*}
    u'(a) = u'(b) = 0.
\end{align*}
Multiplying both sides of $\mathcal{L}u(x) = \lambda u(x)$ by $u(x)$ and integrating, we have
\begin{align*}
    \int_a^b \mathcal{L}u(x) \cdot u(x) ~ dx & =  \int_a^b (p(x)u'(x))' \cdot u(x)  ~ dx                       \\
                                             & =   \Big[p(x) u'(x) u(x)\Big]^b_a - \int_a^b p(x)(u'(x))^2 ~ dx \\
                                             & = - \int_a^b p(x)(u'(x))^2 ~ dx                                 \\
                                             & = \int_a^b \lambda |u(x)|^2 ~ dx.
\end{align*}
That is,
\begin{equation*}
    \lambda = -\int_a^b\frac{ p(x)(u'(x))^2}{|u(x)|^2} ~ dx.
\end{equation*}
Since $p(x) > 0$, $u(x)$ is nontrivial, and the square of reals are non-negative, it follows that $\lambda \leq 0$.

\vspace{5ex}

\noindent{}\textbf{Exercise 2:}
Let $u_\lambda$ be an nontrivial eigenfunction for $\mathcal{L}$ with $\lambda \neq 0$, and let $x_0$ be any critical point on $[A, B]$.

For the sake of contradiction, suppose $u_\lambda(x_0) = 0$. Since $x_0$ is a critical point, it follows that $u'(x_0) = 0$.

Then, we can rewrite the eigenfunction $\mathcal{L}u(x) = \lambda u(x)$ as
\begin{equation}
    u''(x) + \frac{p'(x)}{p(x)}u'(x) - \frac{\lambda}{p(x)}u(x) = 0.
\end{equation}
Notice that this is a second order linear ODE.
So by Picard-Linedelof's Theorem, there must exist a unique solution for each initial condition.

Consider the constant function $v(x) = 0$.
Notice it satisfies our initial conditions $u(x) = u'(x) = 0$, and thus it is a valid solution.
Moreover since $u_\lambda(x_0)$ also satisfies $u(x) = u'(x) = 0$, by uniqueness of solutions, they must be the same.
However, this is impossible since $v(x)$ is the trivial solution and we claimed that $u_\lambda(x_0)$ is a nontrivial solution.
Thus, we have a contradiction.
Therefore, $u_\lambda(x_0) \neq 0$.

Consider equation 1 again. Evaluating at $x_0$:
\begin{align*}
    u''(x_0) + \frac{p'(x_0)}{p(x_0)}u'(x_0) - \frac{\lambda}{p(x_0)}u(x_0) = u''(x_0) - \frac{\lambda}{p(x_0)}u(x_0) = 0.
\end{align*}
Then,
\begin{equation*}
    u''(x_0) = \frac{\lambda}{p(x_0)}u(x_0).
\end{equation*}
From exercise 1, we have that $\lambda \leq 0$ and we let $\lambda \neq 0$, so $\lambda < 0$ must be true.
Since $\lambda < 0$ and $p(x_0) > 0$, it follows that $\frac{\lambda}{p(x_0)} < 0$.
Then, $u''(x_0)$ is $u(x)$ times a negative scalar, and so they must have opposite signs.

Moreover, since we have $u_\lambda(x_0) \neq 0$, $\lambda < 0$, and $p(x_0) > 0$, it follows that $u''(x_0) \neq 0$.

\vspace{5ex}

\noindent{}\textbf{Exercise 3:}
Let $(x_1, x_2) \in [A, B]$ be successive critical points of $u_\lambda$.
Then,
\begin{align*}
    u_\lambda'(x_1) = u_\lambda'(x_2) = 0.
\end{align*}
Define $v_\lambda := u_\lambda |_{[x_1, x_2]}$.
Then $\mathcal{L}v_\lambda(x) = \lambda v_\lambda(x)$ still holds because $[x_1, x_2] \subset [A, B]$.
Since $x_1$ and $x_2$ are critical points, $v_\lambda'(x_1) = 0$ and $v_\lambda'(x_2) = 0$ are both true.
So Neumann boundary conditions are also satisfied.

Now consider the constant function $w(x) = 1 : \forall x \in [x_1, x_2]$.
Notice that $w(x)$ is an eigenfunction of $\mathcal{L}$ with $\lambda = 0$ since
\begin{align*}
    \mathcal{L}w(x) = (p(x) \cdot w'(x))' = 0 = 0 \cdot w(x).
\end{align*}
It also satisfies the Neumann boundary conditions:
\begin{align*}
    w'(x_1) = w'(x_2) = 0.
\end{align*}

Recall that eigenfunctions with different eigenvalues must be orthogonal to each other.
Then, since we have $\lambda \neq 0$:
\begin{align*}
    \int^{x_2}_{x_1} v_\lambda(x) \cdot w(x) ~ dx = \int^{x_2}_{x_1} v_\lambda(x) ~ dx = 0.
\end{align*}

Now suppose, for sake of contradiction, that $v_\lambda(x) \neq 0 : \forall x \in (x_1, x_2)$.
From Exercise 2, we also have that $v_\lambda(x_1), v_\lambda(x_2) \neq 0$. So our bounds expand to
\begin{align*}
    v_\lambda(x) \neq 0 : \forall x \in [x_1, x_2].
\end{align*}

Then $v_\lambda(x)$ is a non-zero continuous function on a bounded interval.
So by Intermediate Value Theorem, $v_\lambda(x)$ never switch signs, and thus we have these two cases.

Case 1:
We have $v_\lambda(x) > 0 : \forall x \in [x_1, x_2]$. Then,
\begin{align*}
    \int^{x_2}_{x_1} v_\lambda(x) > 0.
\end{align*}
This contradicts our earlier findings that $\int^{x_2}_{x_1} v_\lambda(x) = 0$.

Case 2:
We have $v_\lambda(x) < 0 : \forall x \in [x_1, x_2]$. Then,
\begin{align*}
    \int^{x_2}_{x_1} v_\lambda(x) < 0.
\end{align*}
This again contradicts our earlier findings that $\int^{x_2}_{x_1} v_\lambda(x) = 0$.

Since both cases lead to contradictions, it must be that our assumption is false.
Therefore $v_\lambda(x)$ has at least one zero in the interval $(x_1, x_2)$.

\vspace{5ex}

\noindent{}\textbf{Exercise 4:}
Set $\mathcal{L}u(x) = ((1 + x)u'(x))'$, so that $V_{\lambda_{2}}(x) = \frac{1}{4}\frac{1}{(1+x)^2} - \frac{\lambda_2}{1+x}$ for Dirichlet boundary
conditions on $[x_1, x_2]$. Assume $x_1, x_2 > 0$. Differentiating $V_{\lambda_{2}}(x)$, we have
\begin{align*}
    V_{\lambda_{2}}'(x) = -\frac{1}{2}(1+x)^{-3} + \lambda_2 (1+x)^{-2}.
\end{align*}
Since $x_1,x_2 > 0$, $x>0$ in this interval. Moreover, we have that $\lambda_2 < 0$ from the previous exercises.
So $V_{\lambda_{2}}'(x)$ must be strictly negative for all $x \in [x_1, x_2]$.
That is, $V_{\lambda_{2}}(x)$ is monotonically decreasing for all $x \in [x_1, x_2]$.
Then, $V_{\lambda_{2}}(x)$ attains it's maximum at $x_1$ and minimum at $x_2$.
Specifically,
\begin{align*}
    V_{\lambda_{2}}(x_1) = \frac{1}{4}\frac{1}{(1+x_1)^2} - \frac{\lambda_2}{1+x_1} = M^2,
\end{align*}
and
\begin{align*}
    V_{\lambda_{2}}(x_2) = \frac{1}{4}\frac{1}{(1+x_2)^2} - \frac{\lambda_2}{1+x_2} = m^2.
\end{align*}

Recall that $y(x) = \sqrt{p(x)} \cdot u_{\lambda_2}(x)$ where $p(x) >0$ and $u_{\lambda_2}(x)$ is an nontrivial eigenfunction.
So $y(x)$ is a non-trivial solution.
Since we also have that $x_1$ and $x_2$ are successive zeroes, by Theorem 2.1, we have
\begin{align*}
    \frac{\pi}{M} \leq x_2 - x_1 \leq \frac{\pi}{m}.
\end{align*}
We can rearrange this to
\begin{align*}
    m \leq \frac{\pi}{x_2 - x_1} \leq M.
\end{align*}
Squaring everything, we have
\begin{align*}
    m^2 \leq \frac{\pi^2}{(x_2 - x_1)^2} \leq M^2.
\end{align*}
Using our substitutions for $M^2$ and $m^2$, we solve for $\lambda_2$.

For the upper bound:
\begin{align*}
    M^2                                            & \geq \frac{\pi^2}{(x_2 - x_1)^2}                            \\
    \frac{1}{4(1+x_1)^2} - \frac{\lambda_2}{1+x_1} & \geq \frac{\pi^2}{(x_2 - x_1)^2}                            \\
    \frac{\lambda_2}{1+x_1}                        & \leq \frac{1}{4(1+x_1)^2} - \frac{\pi^2}{(x_2 - x_1)^2}     \\
    \lambda_2                                      & \leq \frac{1}{4(1+x_1)}- \frac{\pi^2(1+x_1)}{(x_2 - x_1)^2}
\end{align*}

For the lower bound:
\begin{align*}
    m^2                                            & \leq \frac{\pi^2}{(x_2 - x_1)^2}                            \\
    \frac{1}{4(1+x_2)^2} - \frac{\lambda_2}{1+x_2} & \leq \frac{\pi^2}{(x_2 - x_1)^2}                            \\
    \frac{\lambda_2}{1+x_2}                        & \geq \frac{1}{4(1+x_2)^2} -\frac{\pi^2}{(x_2 - x_1)^2}      \\
    \lambda_2                                      & \geq \frac{1}{4(1+x_2)} -\frac{\pi^2(1+x_2)}{(x_2 - x_1)^2} \\
\end{align*}

Combining, we have
\begin{align*}
    \frac{1}{4(1+x_2)} -\frac{\pi^2(1+x_2)}{(x_2 - x_1)^2}  \leq \lambda_2 \leq \frac{1}{4(1+x_1)}- \frac{\pi^2(1+x_1)}{(x_2 - x_1)^2}.
\end{align*}


\end{document}

